
\section{Decision Problem and Conceptual Framework}
\label{sec:problem}

\subsection{Policy Selection Versus Route Assignment}
\label{sec:problem:levels}

Write $\mathcal{S}$ for the space of operating contexts and $\mathcal{A}$ for the
finite portfolio of routing policies. Two nested decisions are made. The outer
one, which we call \emph{level~1}, observes a context $s \in \mathcal{S}$ and
switches on a single policy $a \in \mathcal{A}$ for the episode. The inner one,
\emph{level~2}, is carried out by that policy: it assigns a path to each guided
vehicle from whatever traffic information it is entitled to read. Learning
occurs only at level~1. The policies themselves are fixed constructions taken
from the literature and are not altered, tuned or combined by the selector.

The separation matters for what can be claimed. Because level~2 is unchanged, an
improvement attributable to level~1 is an improvement in \emph{choosing}, not in
routing, and it is bounded above by the best that the fixed portfolio can do. It
also fixes the unit of analysis: a context is one instance, not one vehicle and
not one replication, and every replication of a context moves together through
every evaluation protocol.

A context contains only quantities defined before the policy is activated. Here
it is the six-tuple of corridor demand, effective green ratio, guidance
penetration, information lag, presence of a disruption and alternative-corridor
capacity. No realised travel time, queue, flow or emission enters it.

\subsection{Context Space and Policy Portfolio}
\label{sec:problem:space}

The context space is the full factorial of the six factors, and the portfolio
contains four established policies, described in Sections~\ref{sec:method:net}
and~\ref{sec:method:pol}. Two properties of the portfolio are treated as
empirical questions rather than assumptions: whether its members are
behaviourally distinct in this environment, and whether their outcome
differences survive replication noise. Both are reported.

\subsection{Decision Cost and Regret}
\label{sec:problem:regret}

Let $C(s,a)$ be the decision cost of policy $a$ in context $s$, measured on the
primary criterion and averaged over the replications of that context. The
primary criterion is the system mean journey time in seconds per vehicle;
Section~\ref{sec:method:crit} states why, and the other two criteria are carried
alongside. Costs are reported in seconds throughout, so no normalisation
convention can affect them.

The per-context best and its cost are
\begin{equation}
a^{\mathrm{VB}}(s) \in \arg\min_{a \in \mathcal{A}} C(s,a),
\qquad
C^{\mathrm{VB}}(s) = \min_{a \in \mathcal{A}} C(s,a).
\label{eq:vbs}
\end{equation}
Selecting $a^{\mathrm{VB}}(s)$ everywhere defines the \emph{virtual best}
reference. It uses the realised performance of every candidate and is therefore
a bound, not a deployable method. For a selection rule
$\pi : \mathcal{S} \to \mathcal{A}$ the \emph{decision regret} in context $s$ is
\begin{equation}
R(s,\pi) \;=\; C\bigl(s,\pi(s)\bigr) - C^{\mathrm{VB}}(s) \;\ge\; 0 .
\label{eq:regret}
\end{equation}
Regret measures the cost of the decision, in seconds per vehicle. It is not a
prediction error, and the two are not monotonically related.

The strongest deployable non-adaptive rule is the \emph{single best} fixed
policy
\begin{equation}
a^{\mathrm{SB}} \;=\; \arg\min_{a \in \mathcal{A}}\;
\frac{1}{|\mathcal{S}|}\sum_{s \in \mathcal{S}} C(s,a),
\label{eq:sbs}
\end{equation}
and the quantity that adaptation competes for is the \emph{decision headroom}
\begin{equation}
H \;=\; \frac{1}{|\mathcal{S}|}\sum_{s \in \mathcal{S}}
\Bigl( C\bigl(s,a^{\mathrm{SB}}\bigr) - C^{\mathrm{VB}}(s) \Bigr),
\qquad
G(\pi) \;=\; 1 - \frac{\overline{R(\pi)}}{H},
\label{eq:headroom}
\end{equation}
where $G(\pi)$ is the share of the headroom a rule recovers: $G=1$ at the
virtual best, $G=0$ at the single best fixed policy and $G<0$ for a rule that is
worse than doing nothing adaptive. In every held-out protocol,
$a^{\mathrm{SB}}$ is identified from training contexts only, so the baseline is
itself a fitted quantity and not an oracle.

\subsection{Complementarity, Decision Value and Deployability}
\label{sec:problem:three}

The framework separates three properties that are often treated as one.

\paragraph{Complementarity} Different policies are best in different contexts:
$a^{\mathrm{VB}}(s)$ is not constant over $\mathcal{S}$. This is a property of
the portfolio and the environment alone. It can be asserted on a scalar
criterion or on the criterion vector, and the two can disagree: a policy may be
non-dominated in many contexts while never minimising any single criterion.
Complementarity says nothing about magnitude.

\paragraph{Decision value} The cost carried by that variation, $H$ in
Eq.~\eqref{eq:headroom}. The two properties come apart whenever the fixed policy
is merely \emph{narrowly} beaten in the contexts it loses: many distinct winners
can then coexist with a vanishing $H$. Complementarity is thus a precondition
for decision value and no more than that, and the two are not even commensurable
--- one counts contexts, the other prices seconds per vehicle.

\paragraph{Deployability} Whether any rule can actually collect $H$ using only
what is knowable before the policy is switched on, in conditions it was not
fitted to, and whether it keeps doing so once those conditions drift. Formally
this is $G(\pi)$ measured out of sample, together with how $\pi$ behaves when the
context distribution is deliberately moved. Decision value does not buy this:
interpolating within a sampled region is a different act from extrapolating
beyond it, and the second can cost more than it returns.

A fourth quantity is needed to interpret the other three. Because every context
is replicated, each has a \emph{replication noise band}: the dispersion of the
paired seed differences between policies. A headroom smaller than that band is
not merely uncertain, it is unobservable at the achieved sample size, and the
distinction between a small real effect and an unresolved one must be reported
rather than absorbed into an average. Section~\ref{sec:method:stats} states the
criterion used, and Section~\ref{sec:value} reports headroom against noise
context by context.
